\documentclass[12pt,a4paper]{article}
\usepackage[utf8]{inputenc}
\usepackage[T1]{fontenc}
\usepackage{amsmath,amssymb}
\usepackage{graphicx}
\usepackage{booktabs}
\usepackage{longtable}
\usepackage[margin=1in]{geometry}
\usepackage{hyperref}
\usepackage{xcolor}

\hypersetup{colorlinks=true, linkcolor=blue!60!black, citecolor=blue!60!black, urlcolor=blue!60!black}

\title{Social Structure Without Social Mechanisms:\\By-Product Mutualism in Energy-Constrained Multi-Agent Systems}

\author{Tristan Stoltz\\Luminous Dynamics\\\texttt{tristan.stoltz@evolvingresonantcocreationism.com}}

\date{2026}

\begin{document}
\maketitle

\begin{abstract}
We report a systematic investigation into the mechanisms underlying emergent cooperation in a thermodynamic multi-agent physics engine. Across 52 findings and 62 experiments, we demonstrate that social structure---spatial clustering, cooperation events, group formation, and survival patterns---emerges as \textbf{by-product mutualism} from individual resource-seeking, not from social mechanisms. Six distinct movement controllers (free-energy-gradient, well-seeking, greedy, partner-seeking, random walk, and stationary) all produce cooperation when energy wells are present ($p < 0.002$, controller-invariant). Disabling mutual energy regeneration has zero effect on well-seeking agents ($d = 0.000$). Scarcity (250--2500\,J), volatility (seasonal, random, migrating), invisibility, threshold access, and resonance-gating all fail to make the social gradient indispensable. The social component of the behavioral gradient is a metabolic liability in every tested environment: it diverts agents from resource access, reducing survival by 1.5--9.4 agents relative to pure well-seeking.

We argue this constitutes a computational proof of the by-product mutualism hypothesis: in environments with locatable spatial resource attractors, the geometry and thermodynamic distribution of resources determines social structure, not the internal algorithms of agents. Social seeking becomes marginally beneficial only under migrating resources ($d = 0.35$) and potentially under resonance-only conditions (experiment pending).

All experiments use 20-seed replication with Holm--Bonferroni correction. The engine is open-source (AGPL-3.0) with bitwise-verified deterministic reproduction.
\end{abstract}

\textbf{Keywords:} by-product mutualism, emergent cooperation, thermodynamic enforcement, multi-agent systems, resource geometry, free energy principle

\section{Introduction}

The emergence of cooperation in multi-agent systems is typically explained through game-theoretic mechanisms: reciprocal altruism \cite{trivers1971}, kin selection \cite{hamilton1964}, reputation \cite{nowak2005}, or institutional enforcement \cite{ostrom1990}. These explanations share an assumption: cooperation requires a \emph{mechanism}---a strategy, a drive, or a rule that causes agents to cooperate.

We present evidence against this assumption. Using Symtropy, an open-source rigid body physics engine with thermodynamic enforcement, we show that social structure---clustering, cooperation events, and group formation---emerges from \emph{individual resource-seeking} in environments with shared spatial resource attractors. No social mechanism is required.

Our investigation proceeded in three phases. First (Findings 1--45), we documented extensive cooperation under a Free Energy Principle (FEP) behavioral gradient and initially attributed it to the gradient's social component. Second (Finding 46), we replaced the FEP gradient with five alternative controllers and found that \emph{all} produce cooperation. Third (Findings 47--52), we systematically tested whether any environmental modification makes the social gradient indispensable. None did.

The result is a computational demonstration of \textbf{by-product mutualism} \cite{connor1995}: agents act in individual self-interest to access resources, and mutual benefits arise because the environment forces them into spatial proximity. This aligns with ecological theory on by-product mutualism and producer-scrounger dynamics \cite{giraldeau2000} but extends it to thermodynamic multi-agent systems with strict energy conservation.

\section{Engine and Methods}

\subsection{Symtropy Engine}

Symtropy is a rigid body physics engine (Rust, 2D--9D, LBVH broadphase, GJK/EPA narrowphase) where agents maintain energy budgets tracked in joules. Each agent pays a maintenance cost per tick ($0.15$--$0.25$\,J/tick) and regenerates energy from three possible sources: energy wells (spatial resource points), harmony resonance (proximity-based mutual benefit with compatible agents), and ambient recovery. When energy reaches zero, the agent collapses.

Agents carry an 8-dimensional harmony profile. When two agents are within the harmony range (40 units) and their profile resonance (cosine similarity) exceeds 0.5, both receive energy regeneration at rate $r_{\text{regen}} \times (R - 0.5) \times 2$\,J/tick, where $R$ is resonance and $r_{\text{regen}} = 0.06$--$0.12$\,J/tick.

\subsection{Controllers}

Six movement controllers were tested, sharing identical thermodynamics:

\begin{enumerate}
\item \textbf{FEP\_GRADIENT}: Handcrafted free-energy-style gradient with four weighted components---seek resonant partners, seek energy wells, flee danger, explore.
\item \textbf{WELL\_ONLY}: Navigate to nearest visible energy well; ignore other agents.
\item \textbf{GREEDY}: 8-direction lookahead; choose direction maximizing next-tick energy gain.
\item \textbf{PARTNER\_ONLY}: Navigate to nearest resonant partner; ignore wells.
\item \textbf{RANDOM}: Uniform random direction each tick.
\item \textbf{STATIONARY}: Do not move.
\end{enumerate}

All controllers share identical energy budgets, maintenance costs, resonance mechanics, and collision physics. Only the movement rule differs. Resonance regeneration is passive---it occurs whenever compatible agents are proximate, regardless of whether they sought each other.

\subsection{Statistical Methods}

All experiments use 20 independent seeds with deterministic LCG random number generation (bitwise-verified by lock-in test). Comparisons use Mann--Whitney $U$ (nonparametric, no normality assumption) with Cohen's $d$ for effect size. Multiple comparisons are corrected by Holm--Bonferroni at $\alpha = 0.05$.

\section{Results}

\subsection{Phase I: Cooperation Appears Universal (F1--45)}

Under the FEP gradient with energy scarcity and finite wells, agents consistently clustered, cooperated, and formed groups. A six-condition ablation (F2) identified the FEP gradient as the clustering mechanism and resonance as the survival mechanism. A metric independence test (F3--4) showed IIT-style $\Phi$, Shannon entropy, constant scalar, and random values all produce identical clustering ($d < 0.1$); only $\Phi = 0$ differed.

We initially interpreted these results as evidence that cooperation is thermodynamically inevitable under scarcity with social range and mutual benefit.

\subsection{Phase II: Controller Replacement Falsifies the Mechanism (F46)}

We replaced the FEP gradient with five alternative controllers (Table~\ref{tab:f46}).

\begin{table}[h]
\centering
\caption{Finding 46: All controllers produce cooperation.}
\label{tab:f46}
\begin{tabular}{lrrrr}
\toprule
Controller & Alive/20 & Cooperation & Clustering & Energy \\
\midrule
FEP\_GRADIENT & 10.8 & 1,081K & 6.37 & 18.2\,J \\
WELL\_ONLY & \textbf{20.0} & \textbf{1,222K} & 1.99 & 200.0\,J \\
PARTNER\_ONLY & 10.5 & 589K & 2.07 & 0.2\,J \\
GREEDY & \textbf{20.0} & \textbf{1,475K} & 3.13 & 200.0\,J \\
RANDOM & 20.0 & 474K & 16.44 & 107.3\,J \\
STATIONARY & 20.0 & 521K & 12.50 & 140.0\,J \\
\bottomrule
\end{tabular}
\end{table}

All six controllers produce substantial cooperation ($> 470$K events). WELL\_ONLY and GREEDY \emph{outperform} the FEP gradient in both survival (20.0 vs 10.8) and cooperation events (1,222K and 1,475K vs 1,081K). The FEP gradient produces the \emph{worst} survival of any controller because its social component diverts agents from wells.

\subsection{Phase III: Systematic Environmental Tests (F47--52)}

\begin{table}[h]
\centering
\caption{Findings 47--52: No environmental modification makes the FEP social gradient indispensable.}
\label{tab:env}
\begin{tabular}{llrrl}
\toprule
Finding & Environment & WELL$\pm$REGEN & FEP$\pm$REGEN & Cooperation necessary? \\
\midrule
F47 & Standard & 20.0 / 20.0 & 11.7 / 3.9 & No ($d = 0.000$ for WELL) \\
F48 & Scarce (250\,J) & 20.0 / 20.0 & --- & No (all capacities) \\
F49 & Hidden wells & 20.0 / --- & 16.8 / --- & No ($d = 1.39$, WELL wins) \\
F50 & Migrating & 9.3 / 6.8 & 9.9 / --- & \textbf{Marginal} ($d = 0.35$) \\
F51 & Threshold ($\geq 4$) & 20.0 / --- & 10.8 / --- & No (natural co-location) \\
F52 & Resonance-gated & 9.3 / --- & 10.9 / --- & No ($d = -0.20$, n.s.) \\
\bottomrule
\end{tabular}
\end{table}

In every environmental condition, the WELL\_ONLY controller matches or exceeds FEP survival. The sole marginal exception is migrating resources (F50), where the FEP gradient converges with WELL\_ONLY ($d = 0.35$, small effect) as both controllers must search for relocated wells.

\section{Discussion}

\subsection{By-Product Mutualism, Not Emergent Cooperation}

Our results demonstrate that the cooperation observed in Findings 1--45 was \textbf{by-product mutualism}---a free benefit of resource co-location, not a product of social mechanisms. In behavioral ecology, by-product mutualism occurs when an organism's selfish actions incidentally benefit others nearby \cite{connor1995}. Our engine instantiates this precisely: agents seeking wells co-locate at shared resources, and proximity-based resonance regeneration provides incidental mutual benefit.

The critical evidence is Finding 47: WELL\_ONLY agents survive identically with and without resonance regeneration ($d = 0.000$). The resonance system provides 95\% of total energy when active but is completely dispensable. Wells alone sustain agents.

\subsection{The FEP Gradient as Metabolic Liability}

The FEP gradient's social component---seeking resonant partners---is counterproductive in every tested environment. It pulls agents away from life-sustaining wells to pursue social contact, creating an energy deficit that resonance then partially compensates. Without resonance (FEP+NO\_REGEN), agents crash to 3.9/20 survival (F47). The social drive creates the problem that cooperation then solves.

This does not invalidate the Free Energy Principle as a theoretical framework. It demonstrates that in environments with locatable spatial resource attractors, the FEP's social component is an \emph{informational} drive that actively competes with basic metabolic survival. The gradient optimizes for prediction error reduction rather than energy efficiency.

\subsection{Resource Geometry Determines Social Structure}

The macroscopic structure of the agent population---clustering, cooperation rate, group size, inequality---is determined by the spatial distribution and accessibility of resources, not by agent behavior. Every controller converges on the same wells, producing the same co-location, the same passive cooperation, and the same social structure.

This is a computational demonstration that \textbf{resource geometry determines the geometry of societies}. The implication extends beyond agent-based models: in any system where agents compete for localized resources, the resource distribution, not the agents' cognitive architecture, is the primary determinant of social organization.

\subsection{When Does Social Seeking Matter?}

The one regime where social seeking shows marginal benefit is resource migration (F50). When wells physically relocate, agents stranded at depleted sites benefit from following other agents who are moving toward new resources. This maps to the producer-scrounger dynamic \cite{giraldeau2000}: the social gradient enables \emph{information scrounging}, not cooperation.

We hypothesize that removing wells entirely---making resonance the sole energy source---would create the regime where social seeking is indispensable (experiment pending). In such a ``resonance-only'' environment, agents cannot survive without finding compatible partners, and the social gradient would transition from liability to necessity.

\subsection{Implications for Prior Findings}

Findings 1--45 remain valid as observational data. Clustering occurs. Phase transitions exist. The Bowling Alone pattern emerges under volitional choice. Adversaries die faster. Equal starts produce maximum inequality. These phenomena are all real---but their cause is resource geometry, not cooperation mechanisms.

The reinterpretation is: ``the geometry of energy wells, not the social drives of agents, produces all observed social structure.'' This is a stronger and more parsimonious explanation than the original.

\subsection{Limitations}

The engine is a stylized model with hand-designed compatibility functions and movement rules. We do not claim these results explain biological or human cooperation. The by-product mutualism result is specific to environments with locatable spatial resource attractors; other environments (resonance-only, large-arena, mobile resources) may produce different dynamics. Statistical power for null results (20 seeds) may miss small effects ($d < 0.4$).

\section{Conclusion}

Social structure requires no social mechanism. In energy-constrained multi-agent systems with shared spatial resources, individual resource-seeking is sufficient to produce clustering, cooperation, and group formation. The behavioral gradient that causes agents to seek social contact is not only unnecessary---it actively reduces survival by diverting agents from resource access.

The cooperation we observed across 45 initial findings was by-product mutualism: a passive consequence of resource co-location, not an emergent property of social dynamics. This result strips the anthropocentric assumption that cooperation requires intent, strategy, or a ``social drive.'' Physics does not cooperate. It clusters at the wells.

\begin{thebibliography}{20}
\bibitem{connor1995} Connor, R.C. (1995). The benefits of mutualism: a conceptual framework. \textit{Biological Reviews}, 70(3), 427--457.
\bibitem{giraldeau2000} Giraldeau, L.-A. \& Caraco, T. (2000). \textit{Social Foraging Theory}. Princeton UP.
\bibitem{hamilton1964} Hamilton, W.D. (1964). The genetical evolution of social behaviour. \textit{J.~Theoretical Biology}, 7(1), 1--16.
\bibitem{nowak2005} Nowak, M.A. \& Sigmund, K. (2005). Evolution of indirect reciprocity. \textit{Nature}, 437, 1291--1298.
\bibitem{ostrom1990} Ostrom, E. (1990). \textit{Governing the Commons}. Cambridge UP.
\bibitem{tononi2004} Tononi, G. (2004). An information integration theory of consciousness. \textit{BMC Neuroscience}, 5(1), 42.
\bibitem{trivers1971} Trivers, R.L. (1971). The evolution of reciprocal altruism. \textit{Quarterly Review of Biology}, 46(1), 35--57.
\bibitem{friston2010} Friston, K. (2010). The free-energy principle: a unified brain theory? \textit{Nature Reviews Neuroscience}, 11(2), 127--138.
\bibitem{prigogine1977} Prigogine, I. (1977). \textit{Self-Organization in Nonequilibrium Systems}. Wiley.
\end{thebibliography}

\end{document}
